El programa se modela para los \textbf{56 meses} contractuales con las paridades del Formulario E-24: 1 UF = 40.000 CLP; 1 USD = 900 CLP; 1 EUR = 1.000 CLP. Se distinguen inversión inicial (CAPEX) y operación continua (OPEX). Los valores siguientes se transcriben de la tabla 3.2 del informe maestro.

\subsection{Costo integral por partida y período contractual}
\begin{center}
\footnotesize\setlength{\tabcolsep}{3pt}\renewcommand{\arraystretch}{1.4}
\begin{tabularx}{\linewidth}{@{}Xrrrrrr@{}}
\toprule
\textbf{Partida (UF netas)} & \textbf{M1--12} & \textbf{M13--20} & \textbf{M21--36} & \textbf{M37--48} & \textbf{M49--56} & \textbf{Total}\\\midrule
CAPEX: ingeniería, cifrado, EIPD y DPA & 375,0 & 0,0 & 0,0 & 0,0 & 0,0 & 375,0\\
CAPEX: certificación inicial ISO 27001 & 0,0 & 387,5 & 0,0 & 0,0 & 0,0 & 387,5\\
DPO & 432,0 & 288,0 & 576,0 & 432,0 & 288,0 & 2.016,0\\
CISO & 576,0 & 384,0 & 768,0 & 576,0 & 384,0 & 2.688,0\\
Soporte QA y Legal & 27,0 & 178,0 & 280,0 & 240,0 & 166,0 & 891,0\\
CISO Assistant Pro Cloud & 60,0 & 40,0 & 80,0 & 60,0 & 40,0 & 280,0\\
Azure Key Vault & 108,0 & 123,5 & 246,96 & 185,2 & 123,5 & 787,16\\
Vigilancia ISO 27001 & 0,0 & 0,0 & 112,5 & 112,5 & 112,5 & 337,5\\
Póliza de ciberriesgo & 90,0 & 60,0 & 120,0 & 90,0 & 60,0 & 420,0\\
Prevención de infracciones y RAT & 15,0 & 70,0 & 85,0 & 85,0 & 55,0 & 310,0\\\midrule
\rowcolor{pale}\textbf{Total} & \textbf{1.683,0} & \textbf{1.531,0} & \textbf{2.268,46} & \textbf{1.780,7} & \textbf{1.229,0} & \textbf{8.492,16}\\
\bottomrule
\end{tabularx}
\end{center}
{\footnotesize\textbf{Tabla 1.} Flujo nominal transcrito del maestro. Fechas de precios declaradas en la fuente: 16/09/2026; certificación: 28/08/2026. La ficha de metadatos y las salvedades de conciliación se encuentran en el anexo A.}

\subsection{Resultados del modelo}
El presupuesto maestro suma \textbf{8.492,16 UF netas} (339.686.400 CLP; aproximadamente USD 377.429). Aplicando la tasa contractual del 0,9\% mensual del Formulario E-24, el documento informa un valor actual del costo de \textbf{6.651,24 UF}.

En los primeros 20 meses se concentra un desembolso de \textbf{3.214,0 UF} (37,8\%). La fase operativa comprende los meses 21 a 56 y el maestro informa un gasto promedio de \textbf{146,6 UF/mes}. Los tramos M21--36 y M49--56 abarcan 16 y 8 meses, respectivamente; no son años completos.

\subsection{Puente de conciliación presupuestaria}
\[
\mathrm{TCO}=\underbrace{4.855,16\ \mathrm{UF}}_{\text{obligaciones directas}}
+\underbrace{3.637,0\ \mathrm{UF}}_{\text{gobernanza y soporte}}
=\boxed{8.492,16\ \mathrm{UF}}.
\]
Esta inversión representa aproximadamente un \textbf{3,95\%} del presupuesto total estimado para la licitación del Caso 10 (aproximadamente 215.000 UF a 56 meses), según la relación consignada en el maestro.

\subsection{Vínculo con la propuesta técnico-económica}
Cada control se vincula con un paquete de trabajo de la EDT y con los hitos del Formulario E-25. Los perfiles CISO, DPO, QA y Asesor Legal TIC se valorizan mediante las referencias del Formulario E-26. El presupuesto permite identificar el costo de implantar controles y el costo recurrente de sostenerlos durante el contrato.

\nota{\textbf{Alcance de las cifras.} Los valores corresponden al modelo del informe maestro, no a nuevas cotizaciones obtenidas durante la diagramación. El TCO y el VAN reproducibles se revisan mediante el código auxiliar; las discrepancias internas de la fuente se explicitan en el anexo A.}
